\documentclass[10pt,a4paper]{article}
\usepackage[T1]{fontenc}
\usepackage[utf8]{inputenc}
\usepackage[margin=18mm]{geometry}
\PassOptionsToPackage{hyphens}{url}
\usepackage{textcomp,url,microtype,enumitem}
\usepackage[hidelinks]{hyperref}
\setlength{\parindent}{0pt}
\setlength{\parskip}{4pt}
\setlength{\emergencystretch}{2em}
\pagestyle{empty}
\hypersetup{pdfauthor={Hongsheng Wang, Zihan Zeng, Fei Wu, Feng Lin and Xingxu Huang}}
\begin{document}
26 November 2026\par
Prof. Quan Zou (Editor in Chief) and Dr. Zhenxia Chen (Deputy Editor)\par
\textit{Briefings in Bioinformatics}\par
Manuscript ID: BIB-26-2683 (Major Revision)\par
Article type: Problem solving protocol\par
Title: \textit{Separating reconstruction, generation and antibacterial activity in antimicrobial peptide design}\par
Authors: Hongsheng Wang, Zihan Zeng, Fei Wu, Feng Lin and Xingxu Huang

Dear Prof. Zou and Dr. Chen,

We thank you and the three reviewers for the constructive, thorough evaluation of our manuscript (BIB-26-2683). We are pleased to submit our substantially strengthened Major Revision, addressing every major and minor concern raised by all three reviewers through new multi-seed computational experiments on an 11-fold larger, strictly homology-isolated independent test benchmark, complete wet-lab protocol confirmation, expanded sequence-quality and memorization audits, streamlined prose and explicit overlap disclosure.

Key additions in this revision include:
\begin{itemize}[leftmargin=1.4em,itemsep=2pt,topsep=3pt]
\item \textbf{Expanded homology-isolated benchmark with untouched independent test set:} Alongside our original 279/41 pilot cohort, we built and froze a 3,565-peptide DBAASP benchmark (3,120 training, 223 validation, 222 untouched independent test peptides; 5,083 test residues) partitioned by $\geq$70\% global Needleman--Wunsch connected components with zero cross-split or pilot homology $\geq$70\%.
\item \textbf{Multi-seed validation, large-sample generation and statistical testing:} All core experiments now span 3--5 independent seeds with 128--256 generated sequences per terminal condition, 95\% confidence intervals and paired significance tests, alongside $256\times 5$ large-sample replays of all four frozen initial-submission checkpoints.
\item \textbf{KL-$\beta$, 21-fold MSE scaling, padding ablation and capacity-matched comparators:} We directly tested KL weights $\beta\in\{0.01,0.1,1.0\}$, 21-fold MSE loss rescaling with tuned learning rate ($2\times 10^{-3}$) and padding masking, and added three closed-form statistical baselines (Unigram, Bigram and Trigram Markov), capacity-matched autoregressive generators (LSTM-matched 1.34M, Transformer-LM 1.80M) and matched non-diffusion GRU-VAEs (1.99M). Notably, tuning lifts independent-test reconstruction to 96.05\% (CE21) and 94.24\% (tuned MSE21), yet both VAE--diffusion models still trail capacity-matched autoregressive generators by $\sim$20-fold on FREE MMD$^2$---strengthening our central conclusion.
\item \textbf{Metric redundancy, amphipathicity and near-duplicate memorization audits:} PCA auditing shows the six low-order metrics span 1.54 effective dimensions (PC1+PC2 = 95.3\%), while tripeptide JS, $\alpha$-helical hydrophobic moment ($\mu_H$) and $\geq$70\% nearest-neighbor alignment identity separate neural generators from Unigram sampling ($p\leq 0.028$) and reveal up to 34.1\% hidden near-homolog collapse despite 0\% exact copies.
\item \textbf{Confirmed broth microdilution protocol, 1,000$\rightarrow$30$\rightarrow$27 provenance funnel and AMPGen overlap matrix:} We now document the complete confirmed MHB broth microdilution protocol (1\% DMSO, $5\times 10^5$ CFU/mL, 37~$^\circ$C, 16--20 h), justify the OD$_{600}<0.10$ cutoff via a six-rule sensitivity audit (confirming 6/27 \textit{S. aureus} hit invariance across five rules and 3.09--12.87 \textmu{}M potency), explicitly emphasize the absence of haemolysis/cytotoxicity measurements, and tabulate zero candidate sequence overlap (0/27) with AMPGen.
\item \textbf{Requested references and editorial compliance:} We integrated both requested references (PMID: 34730875; PMID: 33848577) plus three verified peptide-analysis citations, highlighted all manuscript changes in blue, included alternative text under every figure legend, and pasted our complete point-by-point response directly into the ScholarOne response text box.
\end{itemize}

Sincerely,\par
\textbf{Hongsheng Wang}\par
Corresponding author\par
College of Computer Science and Technology, Zhejiang University\par
\href{mailto:wanghongsheng@zju.edu.cn}{wanghongsheng@zju.edu.cn}
\end{document}
